\vspace{-3.2em}
\begin{abstract}
\footnotesize
\linespread{0.96}\selectfont
\noindent
Quanto do mercado de trabalho brasileiro está efetivamente sujeito à inteligência artificial --- e quem são os trabalhadores expostos? A partir de um modelo de designação de tarefas com adoção corporativa e individual de IA sob informalidade, este trabalho avalia as previsões teóricas nos microdados da Pesquisa Nacional por Amostra de Domicílios Contínua (PNAD Contínua/IBGE, referente ao primeiro trimestre de 2026), combinados deterministamente ao índice de \citet{eloundou2024gpts}. A exposição concentra-se no núcleo formal e escolarizado: na base, 35\% da força de trabalho (35{,}7 milhões de pessoas) atua em ocupações de exposição quase nula ($\theta \le 2$). Em modelos de Mínimos Quadrados Ponderados (WLS) no nível individual com controles mincerianos completos e inferência agrupada em 122 ocupações ($N = 227.629$), cada ano adicional de estudo associa-se a 0{,}23 ponto a mais de exposição à IA (erro-padrão 0{,}03). A associação bruta negativa entre exposição e informalidade ($-6{,}23$ p.p.) atenua-se em 37{,}2\% após o controle educacional, retendo 62{,}8\% de efeito direto decorrente do capital organizacional corporativo. Nas 27 Unidades da Federação, o gradiente educacional intensifica-se com a profundidade do setor formal ($\text{interação} = 0{,}28$), operando um limiar de reversão em $e^* = 11{,}14$ anos de estudo. A informalidade funciona como o principal amortecedor da exposição à IA, indicando que as políticas de capacitação e transição tecnológica devem priorizar o estrato médio-formalizado.\par
\vspace{0.25em}
\noindent \textbf{Palavras-chave:} Inteligência Artificial, Mercado de Trabalho, Informalidade, Capital Organizacional, Pareamento Ocupacional, PNAD Contínua.\par
\vspace{0.1em}
\noindent \textbf{Classificação JEL:} J23, J24, O33, O17, J46.\par
\vspace{0.25em}
\noindent \rule{\linewidth}{0.4pt}\par
\vspace{0.25em}
\noindent \textbf{Abstract:}
How much of the Brazilian labor market is effectively exposed to artificial intelligence, and who are the exposed workers? Developing a task assignment model with corporate and individual AI adoption under contractual informality, this paper tests theoretical predictions on microdata from the Brazilian Continuous National Household Sample Survey (PNAD Contínua/IBGE, 2026Q1) deterministically linked to the \citet{eloundou2024gpts} occupational exposure index. Technological exposure is concentrated in the formal, highly-educated core: at the bottom, 35\% of the workforce (35.7 million workers) works in near-zero exposure occupations ($\theta \le 2$). Individual-level Weighted Least Squares (WLS) regressions with full Mincerian controls and standard errors clustered across 122 occupations ($N = 227,629$) reveal that each additional year of schooling is associated with a 0.23-point increase in AI exposure (SE 0.03). The gross negative association between exposure and informality ($-6.23$ p.p.) attenuates by 37.2\% once individual education is accounted for, retaining a 62.8\% direct organizational governance channel. Across all 27 Brazilian states, the educational gradient steepens with formal sector maturity ($\text{interaction} = 0.28$), with an inflection threshold at $e^* = 11.14$ years of schooling. Contractual informality serves as the primary structural buffer against AI exposure, indicating that technological adaptation and upskilling programs should target the mid-skilled formal stratum.\par
\vspace{0.25em}
\noindent \textbf{Keywords:} Artificial Intelligence, Labor Markets, Informality, Organizational Capital, Occupational Sorting, PNAD Contínua.\par
\vspace{0.1em}
\noindent \textbf{JEL Codes:} J23, J24, O33, O17, J46.
\end{abstract}
\newpage

